\subsection{Общий вид критерия согласия}

\textbf{Критерий} для проверки выдвинутых гипотез --- это функция вида:
\[
    \delta(\vec{X})=
    \begin{cases}
        H_1, &\vec{X}\notin S\\
        H_2, &\vec{X}\in S
    \end{cases},
\]

где \( S \) --- \textit{критическая область} --- область в множестве всех возможных значений выборки, где принимается альтернативная гипотеза.

Чтобы описать критическую область \( S \) аналитически, введём \textit{функцию отклонения} \( \rho \) --- статистику со свойствами:

\begin{enumerate}
    \item Если верна гипотеза \( H_1 \), то для известного распределения \( \mathcal{G} \) верно:
    \[
        \rho(\vec{X})\sim\mathcal{G}\qquad\text{или}\qquad\rho(\vec{X})\xrightarrow[n\to\infty]{}\eta\sim\mathcal{G}
    \]
    \item Если верна гипотеза \( H_2 \), то:
    \[
        |\rho(\vec{X})|\xrightarrow[n\to\infty]{p}\infty
    \]
\end{enumerate}

Зададим критическую область следующим образом:
\[
    S=\{\vec{X}\mid\rho(\vec{X})\geq C\}
\]

Константа \( C \) определяется из уровня значимости \( \varepsilon=0.05 \):
\[
    \varepsilon=P\{|\eta|\geq C\}
\]

Описанный вид критерия для проверки сформулированных гипотез называется \textit{критерием согласия}.